\chapter{Background: Text Representation and Retrieval}
Search and retrieval are the core mechanisms for accessing internship listings once they have been processed. In this chapter, we explore the evolution from traditional lexical matching to modern semantic search techniques, highlighting the trade-offs in accuracy and performance as established in fundamental retrieval theory \cite{manning2008introduction, baeza2011modern, croft2010search}.

\section{Keyword Retrieval}
Keyword-based retrieval, also known as lexical or sparse retrieval, relies on the presence of exact words or tokens in both the query and the document. This approach has been the foundation of information retrieval for decades \cite{manning2008introduction, zobel2006inverted}.

\subsection{Inverted Indexing}
The fundamental data structure for keyword search is the \textit{inverted index}. It maps each unique term (token) to a list of document identifiers (postings) containing that term. This allows for sub-linear search time regardless of the total number of documents \cite{manning2008introduction, zobel2006inverted}.

\subsection{BM25 (Best Matching 25)}
While simple boolean retrieval checks for term presence, ranking models like BM25 provide a sophisticated scoring mechanism based on probabilistic principles \cite{robertson2009probabilistic, jurafsky2023speech, buttcher2016information}. BM25 improves upon the classic TF-IDF (Term Frequency-Inverse Document Frequency) by introducing two key parameters:
\begin{itemize}
    \item \textbf{Term Frequency Saturation ($k_1$):} It limits the contribution of a term that appears many times in a single document, preventing "keyword stuffing" from overly influencing the score.
    \item \textbf{Length Normalization ($b$):} It penalizes longer documents that might contain a term many times simply due to their length, ensuring that concise, relevant postings are not overshadowed.
\end{itemize}
BM25 remains a powerful baseline, especially for domain-specific technical terms where exact matching is critical \cite{robertson2009probabilistic, thakur2021beir}.

\section{Vector Representations}
Vector-based retrieval (dense retrieval) represents text as high-dimensional numerical vectors, where the geometric distance between vectors correlates with the semantic similarity of the underlying text \cite{jurafsky2023speech, luan2021sparse, lin2021pretrained}.

\subsection{Dense Embeddings}
Unlike sparse vectors, dense embeddings are typically 384 to 1536 dimensions long. These vectors are generated by pre-trained neural networks, most notably those based on the Transformer architecture \cite{vaswani2017attention, devlin2019bert}.

\subsection{Sentence-BERT (SBERT) and Bi-Encoders}
Standard BERT models are not naturally suited for large-scale similarity searches because they require "cross-encoding" which is computationally prohibitive. Sentence-BERT introduced the \textit{bi-encoder} architecture \cite{reimers2019sentence}, which has become a standard for efficient dense retrieval alongside frameworks like Dense Passage Retrieval (DPR) \cite{karpukhin2020dense}. It processes the query and the document independently to produce fixed-sized embeddings. This allows for the pre-computation and indexing of document vectors, making real-time semantic search feasible \cite{reimers2019sentence, johnson2019billion}.

\section{Hybrid Retrieval}
In the context of internship listings, neither keyword nor vector search is perfect in isolation. Hybrid retrieval combines the strengths of both, often outperforming single-modality systems in zero-shot or out-of-distribution scenarios \cite{thakur2021beir, luan2021sparse}:
\begin{enumerate}
    \item \textbf{Sparse Retrieval (BM25)} handles exact matches, acronyms, and rare technical terms. Recent advances like SPLADE \cite{formal2021splade} further enhance this by using neural expansions.
    \item \textbf{Dense Retrieval (SBERT)} handles synonyms, paraphrasing, and conceptual queries \cite{reimers2019sentence}.
\end{enumerate}

\subsection{Reciprocal Rank Fusion (RRF)}
A common method to combine these disparate scores is Reciprocal Rank Fusion. It calculates a final score based on the rank of a document in each retrieval system, rather than the raw score values \cite{cormack2009reciprocal}. This method is robust and does not require parameter tuning across different scales of scores.

\section{Efficiency Considerations}
Scaling semantic search to thousands of job postings requires specialized infrastructure to manage high-dimensional vector operations \cite{johnson2019billion, wang2021milvus}.

\subsection{Vector Databases and ANN}
Traditional SQL databases are not optimized for nearest neighbor searches in high-dimensional space. Vector databases utilize \textit{Approximate Nearest Neighbor (ANN)} algorithms to trade a small amount of accuracy for significant speedups \cite{malkov2018efficient, wang2021milvus}.

\subsection{HNSW (Hierarchical Navigable Small World)}
HNSW is one of the most popular ANN algorithms. It constructs a multi-layered graph where the top layers contain long-range edges for fast navigation and the bottom layers contain short-range edges for local precision \cite{malkov2018efficient}. This structure allows for logarithmic search complexity, making it suitable for low-latency production environments.

\subsection{Quantization}
To reduce memory footprint, techniques like \textit{Scalar Quantization} or \textit{Product Quantization (PQ)} are used \cite{jegou2010product, guo2020accelerating}. These compress floating-point vectors into lower-precision representations, significantly reducing the RAM required to store the index while maintaining acceptable retrieval recall \cite{menghani2023efficient}.
